\begin{lemma}
Assume \(\noderef{ParameterHierarchy}(\varepsilon,\varepsilon_1,\varepsilon_2,
n_0)\), \(\beta=\frac12\), \(n>n_0\),
\[
M=\noderef{TwoBiteNaturalReducedVertexCount}(n),\qquad
K=\noderef{TwoBiteNaturalIndependenceScale}(1+\varepsilon,n),
\qquad
p=\noderef{TwoBiteEdgeProbability}\!\left(\frac12,n\right).
\]
Then
\[
\Pr\bigl(\noderef{FiberAndDegreeEvent}\ \text{holds and }
\noderef{TwoBiteSmallClosedPairsEvent}\ \text{fails}\bigr)\le \frac1n.
\]
\end{lemma}
\begin{proof}
Let
\[
\mathbb P_\beta(E)=
\noderef{TwoBiteEventProbability}
  \bigl(n,M,\noderef{TwoBiteEdgeProbability}(\beta,n),E\bigr)
\]
be the parent probability law after rewriting the Lean parameters by the
rounded equalities \(m=M\) and \(k=K\) from the statement.  For every event
\(E\), the hypothesis \(\beta=\frac12\) gives the pointwise equality
\[
\mathbb P_\beta(E)=
\noderef{TwoBiteEventProbability}
  \bigl(n,M,\noderef{TwoBiteEdgeProbability}(\tfrac12,n),E\bigr).
\]
The statement also names
\[
p=\noderef{TwoBiteEdgeProbability}\!\left(\frac12,n\right),
\]
so this is the same finite sample law as
\[
\mathbb P(E)=
\noderef{TwoBiteEventProbability}(n,M,p,E).
\]
Thus it is enough to prove the desired estimate under \(\mathbb P\); rewriting
the edge parameter by \(\beta=\frac12\), and then replacing \(M,K\) by the
Lean parameters \(m,k\), gives the Lean statement exactly.

Write
\[
\mathcal D(\omega):=\noderef{FiberAndDegreeEvent}(\omega,\varepsilon_2)
\]
and
\[
\mathcal S(\omega):=
\noderef{TwoBiteSmallClosedPairsEvent}
  (k=K,\varepsilon,\varepsilon_1)(\omega).
\]
Unfolding \(\noderef{TwoBiteSmallClosedPairsEvent}\), and then
\(\noderef{ClosedPairsBound}\), gives
\[
\mathcal S(\omega)\quad\Longleftrightarrow\quad
\forall I\subseteq\mathrm{Fin}(n),\ |I|=K\Longrightarrow
\left|\noderef{TwoBiteClosedPairsInClass}
  \bigl(\omega,I,\noderef{TwoBiteSmallClass}(\omega,\varepsilon,I)\bigr)
\right|\le\varepsilon_1K^2. \tag{1}
\]
Equivalently, the \(I\)-instance of \(\mathcal S\) fails exactly when
\[
\neg\noderef{ClosedPairsBound}\!\left(
\left|\noderef{TwoBiteClosedPairsInClass}
  \bigl(\omega,I,\noderef{TwoBiteSmallClass}(\omega,\varepsilon,I)\bigr)
\right|,
\varepsilon_1,K\right)
\]
holds.  Since the sample space and the collection of finsets of
\(\mathrm{Fin}(n)\) are finite, the negation of the universal statement in
(1) is the ordinary finite existential: \(\mathcal S(\omega)\) fails iff
there is a finset \(I\subseteq\mathrm{Fin}(n)\) with \(|I|=K\) for which this
same \(\noderef{ClosedPairsBound}\) predicate fails.

Let
\[
\mathcal I_K=\{I\subseteq\mathrm{Fin}(n): |I|=K\}.
\]
Equivalently,
\[
\mathcal S^c
=\left\{\omega:\exists I\in\mathcal I_K,\ 
\neg\noderef{ClosedPairsBound}\!\left(
\left|\noderef{TwoBiteClosedPairsInClass}
  \bigl(\omega,I,\noderef{TwoBiteSmallClass}(\omega,\varepsilon,I)\bigr)
\right|,
\varepsilon_1,K\right)\right\}. \tag{1a}
\]
For \(I\in\mathcal I_K\), define the fixed-set bad event
\[
E_I=\left\{\omega:\mathcal D(\omega)\ \text{and }
\neg\noderef{ClosedPairsBound}\!\left(
\left|\noderef{TwoBiteClosedPairsInClass}
  \bigl(\omega,I,\noderef{TwoBiteSmallClass}(\omega,\varepsilon,I)\bigr)
\right|,
\varepsilon_1,K\right)\right\}.
\]
Then
\[
\mathcal D\cap\mathcal S^c
=\bigcup_{\substack{I\subseteq\mathrm{Fin}(n)\\ |I|=K}} E_I. \tag{2}
\]
Indeed, if \(\omega\in\mathcal D\cap\mathcal S^c\), then \(\mathcal D\) holds
and the negation of the universal quantifier in (1) gives a finset
\(I\subseteq\mathrm{Fin}(n)\) with \(|I|=K\) for which the
\(\noderef{ClosedPairsBound}\) predicate fails.  Hence
\(\omega\in E_I\).  Conversely, if \(\omega\in E_I\) for some such \(I\), then
\(\mathcal D\) holds and the \(I\)-instance of \(\mathcal S\) fails, so
\(\omega\in\mathcal D\cap\mathcal S^c\).  The event \(\mathcal D\) is thus
retained inside every fixed-set event.

If \(K>n\), then there is no finset \(I\subseteq\mathrm{Fin}(n)\) with
\(|I|=K\), since every finset of \(\mathrm{Fin}(n)\) has cardinality at most
\(|\mathrm{Fin}(n)|=n\).  Thus \(\mathcal I_K=\varnothing\).  The universal
condition defining \(\mathcal S\) has no \(K\)-set instance to fail, so
\(\mathcal S^c=\varnothing\), equivalently the union in (2) is empty.  Hence
for every sample \(\omega\) the indicator of
\(\mathcal D\cap\mathcal S^c\) is \(0\).  The finite sum defining its
probability therefore has no contributing samples, and
\[
\mathbb P(\mathcal D\cap\mathcal S^c)=0\le(n:\mathbb R)^{-1}.
\]
The last inequality uses \(n>n_0\) and \(n_0\in\mathbb N\), so \(n\ge1\) and
the reciprocal of \(n\) is nonnegative.
This proves the result in that case.  Hence assume \(K\le n\) from now on.
The \(K\)-subsets of \(\mathrm{Fin}(n)\) are the elements of
\(\mathrm{powersetCard}(K,\mathrm{Fin}(n))\), so
\[
|\mathcal I_K|=\texttt{Nat.choose}(n,K)=\binom{n}{K}. \tag{3}
\]

Fix \(I\in\mathcal I_K\).  The last conjunct of
\(\noderef{ParameterHierarchy}(\varepsilon,\varepsilon_1,\varepsilon_2,n_0)\)
is precisely
\(\noderef{ParameterHierarchyEventualComparisons}
(\varepsilon,\varepsilon_1,\varepsilon_2,n_0)\).  Apply
\(\noderef{SmallClosedPairsFixedSetTail}\) with this package, with the
inequality \(n>n_0\), and with the exact rounded parameter equalities
\[
M=\noderef{TwoBiteNaturalReducedVertexCount}(n),\qquad
K=\noderef{TwoBiteNaturalIndependenceScale}(1+\varepsilon,n),
\qquad
p=\noderef{TwoBiteEdgeProbability}\!\left(\frac12,n\right),
\]
as well as the hypothesis \(|I|=K\).  The child theorem is stated for the law
\[
\noderef{TwoBiteEventProbability}(n,M,p,\cdot),
\qquad
p=\noderef{TwoBiteEdgeProbability}\!\left(\frac12,n\right),
\]
which is the same law \(\mathbb P\) obtained from the parent law after using
\(\beta=\frac12\).  The rounded equalities \(m=M\) and \(k=K\) are exactly the
equalities needed to match the child theorem's \(m\)- and \(k\)-parameters;
there is no additional change of probability space.  Its event is exactly
\(E_I\): both events require
\(\noderef{FiberAndDegreeEvent}(\omega,\varepsilon_2)\), and both require the
negation of
\[
\noderef{ClosedPairsBound}\!\left(
\left|\noderef{TwoBiteClosedPairsInClass}
  \bigl(\omega,I,\noderef{TwoBiteSmallClass}(\omega,\varepsilon,I)\bigr)
\right|,
\varepsilon_1,K\right)
\]
for the same fixed \(I\), the same small class, and the same bound
\(\varepsilon_1K^2\).  Therefore
\[
\mathbb P(E_I)\le
\exp\!\left(
-\frac{\binom{K}{2}_{\mathbb R}^2}
 {(\log n)M n^{8\varepsilon}}
\right). \tag{4}
\]
Here the \(K\) in \(\binom{K}{2}_{\mathbb R}\) is the real coercion of the
natural number \(K\), and \(M\) in the denominator is the real coercion of the
natural number \(M\), exactly as in the Lean statement of
\(\noderef{SmallClosedPairsFixedSetTail}\).

Apply the finite union bound to the union (2).  No disjointness is assumed.
Indeed, by unfolding \(\noderef{TwoBiteEventProbability}\), the probability
of an event \(A\) is the finite sum
\[
\mathbb P(A)=
\sum_{\omega}
\mathbf 1_A(\omega)\,\noderef{TwoBiteSampleWeight}(\omega).
\]
The weights \(\noderef{TwoBiteSampleWeight}(\omega)\) are nonnegative: by
\(\noderef{TwoBiteSampleWeight}\) they are products of the two
\(G(M,p)\)-graph weights and the uniform injection weight, all of which are
probability weights.  For each sample \(\omega\),
\[
\mathbf 1_{\bigcup_{I\in\mathcal I_K}E_I}(\omega)
\le \sum_{I\in\mathcal I_K}\mathbf 1_{E_I}(\omega),
\]
because if the left side is \(1\), then \(\omega\) lies in at least one of the
events \(E_I\), while if the left side is \(0\) the inequality is immediate.
Multiplying this pointwise inequality by the nonnegative sample weight and
summing over the finite sample space gives
\[
\mathbb P\!\left(\bigcup_{I\in\mathcal I_K}E_I\right)
\le \sum_{I\in\mathcal I_K}\mathbb P(E_I).
\]
Using (4) for every \(I\in\mathcal I_K\) and then using the count (3) gives
\[
\mathbb P(\mathcal D\cap\mathcal S^{c})
\le
\sum_{I\in\mathcal I_K}\mathbb P(E_I)
\le
\sum_{I\in\mathcal I_K}
\exp\!\left(
-\frac{\binom{K}{2}_{\mathbb R}^2}
 {(\log n)M n^{8\varepsilon}}
\right)
=
\exp\!\left(
-\frac{\binom{K}{2}_{\mathbb R}^2}
 {(\log n)M n^{8\varepsilon}}
\right)\binom{n}{K}. \tag{5}
\]
The McDiarmid-union conjunct of
\(\noderef{ParameterHierarchyEventualComparisons}\), instantiated at this
same \(n>n_0\), uses the rounded definitions
\[
K=\noderef{TwoBiteNaturalIndependenceScale}(1+\varepsilon,n),
\qquad
M=\noderef{TwoBiteNaturalReducedVertexCount}(n),
\]
and is precisely the inequality
\[
\exp\!\left(
-\frac{\binom{K}{2}_{\mathbb R}^2}
 {(\log n)M n^{8\varepsilon}}
\right)
\binom{n}{K}\le (n:\mathbb R)^{-1}. \tag{6}
\]
Combining (5) and (6) proves
\[
\mathbb P(\mathcal D\cap\mathcal S^{c})\le(n:\mathbb R)^{-1}.
\]
Finally, replacing \(\mathbb P\) by the parent law
\(\mathbb P_\beta\) is justified by \(\beta=\frac12\), as observed at the
start.  This is the claimed probability bound.
\end{proof}
